% Options for packages loaded elsewhere
\PassOptionsToPackage{unicode}{hyperref}
\PassOptionsToPackage{hyphens}{url}
\PassOptionsToPackage{dvipsnames,svgnames,x11names}{xcolor}
%
\documentclass[
  10pt,
]{article}
\usepackage{amsmath,amssymb}
\usepackage{iftex}
\ifPDFTeX
  \usepackage[T1]{fontenc}
  \usepackage[utf8]{inputenc}
  \usepackage{textcomp} % provide euro and other symbols
\else % if luatex or xetex
  \usepackage{fontspec}
  \defaultfontfeatures{Scale=MatchLowercase}
  \defaultfontfeatures[\rmfamily]{Ligatures=TeX,Scale=1}
\fi
\usepackage{lmodern}
\ifPDFTeX\else
  % xetex/luatex font selection
    \setmainfont[]{STIX Two Text}
\fi
% Use upquote if available, for straight quotes in verbatim environments
\IfFileExists{upquote.sty}{\usepackage{upquote}}{}
\IfFileExists{microtype.sty}{% use microtype if available
  \usepackage[]{microtype}
  \UseMicrotypeSet[protrusion]{basicmath} % disable protrusion for tt fonts
}{}
\makeatletter
\@ifundefined{KOMAClassName}{% if non-KOMA class
  \IfFileExists{parskip.sty}{%
    \usepackage{parskip}
  }{% else
    \setlength{\parindent}{0pt}
    \setlength{\parskip}{6pt plus 2pt minus 1pt}}
}{% if KOMA class
  \KOMAoptions{parskip=half}}
\makeatother
\usepackage{xcolor}
\usepackage[margin=1in]{geometry}
\usepackage{longtable,booktabs,array}
\usepackage{calc} % for calculating minipage widths
% Correct order of tables after \paragraph or \subparagraph
\usepackage{etoolbox}
\makeatletter
\patchcmd\longtable{\par}{\if@noskipsec\mbox{}\fi\par}{}{}
\makeatother
% Allow footnotes in longtable head/foot
\IfFileExists{footnotehyper.sty}{\usepackage{footnotehyper}}{\usepackage{footnote}}
\makesavenoteenv{longtable}
\usepackage{graphicx}
\makeatletter
\newsavebox\pandoc@box
\newcommand*\pandocbounded[1]{% scales image to fit in text height/width
  \sbox\pandoc@box{#1}%
  \Gscale@div\@tempa{\textheight}{\dimexpr\ht\pandoc@box+\dp\pandoc@box\relax}%
  \Gscale@div\@tempb{\linewidth}{\wd\pandoc@box}%
  \ifdim\@tempb\p@<\@tempa\p@\let\@tempa\@tempb\fi% select the smaller of both
  \ifdim\@tempa\p@<\p@\scalebox{\@tempa}{\usebox\pandoc@box}%
  \else\usebox{\pandoc@box}%
  \fi%
}
% Set default figure placement to htbp
\def\fps@figure{htbp}
\makeatother
\setlength{\emergencystretch}{3em} % prevent overfull lines
\providecommand{\tightlist}{%
  \setlength{\itemsep}{0pt}\setlength{\parskip}{0pt}}
\setcounter{secnumdepth}{-\maxdimen} % remove section numbering
% Map the Unicode math operators used in prose to LaTeX math commands,
% so they render even though the text font lacks the glyphs.
\usepackage{amssymb,amsmath}
\usepackage{newunicodechar}
\newunicodechar{→}{\ensuremath{\rightarrow}}
\newunicodechar{∂}{\ensuremath{\partial}}
\newunicodechar{∈}{\ensuremath{\in}}
\newunicodechar{∎}{\ensuremath{\blacksquare}}
\newunicodechar{□}{\ensuremath{\Box}}
\newunicodechar{∝}{\ensuremath{\propto}}
\newunicodechar{∞}{\ensuremath{\infty}}
\newunicodechar{∫}{\ensuremath{\int}}
\newunicodechar{≈}{\ensuremath{\approx}}
\newunicodechar{≠}{\ensuremath{\neq}}
\newunicodechar{≡}{\ensuremath{\equiv}}
\newunicodechar{≤}{\ensuremath{\leq}}
\newunicodechar{≥}{\ensuremath{\geq}}
\newunicodechar{≪}{\ensuremath{\ll}}
\newunicodechar{≫}{\ensuremath{\gg}}
\newunicodechar{≲}{\ensuremath{\lesssim}}
\newunicodechar{⊙}{\ensuremath{\odot}}
\newunicodechar{⊗}{\ensuremath{\otimes}}
\newunicodechar{⟹}{\ensuremath{\Longrightarrow}}
\newunicodechar{⟺}{\ensuremath{\Longleftrightarrow}}
\newunicodechar{⇒}{\ensuremath{\Rightarrow}}
\newunicodechar{⇔}{\ensuremath{\Leftrightarrow}}
\newunicodechar{𝓗}{\ensuremath{\mathcal{H}}}
\newunicodechar{√}{\ensuremath{\surd}}
\newunicodechar{∼}{\ensuremath{\sim}}
\newunicodechar{∇}{\ensuremath{\nabla}}
\newunicodechar{↔}{\ensuremath{\leftrightarrow}}
\newunicodechar{⟨}{\ensuremath{\langle}}
\newunicodechar{⟩}{\ensuremath{\rangle}}
% keep figures from floating far from their callouts
\usepackage{float}
\let\origfigure\figure
\let\endorigfigure\endfigure
\usepackage{bookmark}
\IfFileExists{xurl.sty}{\usepackage{xurl}}{} % add URL line breaks if available
\urlstyle{same}
\hypersetup{
  pdftitle={Is the volume-law horizon a postulate or a consequence? A driven non-equilibrium reduction of Structural Entropy Dark Energy},
  pdfauthor={Stilian Pandev},
  colorlinks=true,
  linkcolor={RoyalBlue},
  filecolor={Maroon},
  citecolor={Blue},
  urlcolor={RoyalBlue},
  pdfcreator={LaTeX via pandoc}}

\title{Is the volume-law horizon a postulate or a consequence? A driven
non-equilibrium reduction of Structural Entropy Dark Energy}
\author{Stilian Pandev}
\date{30 June 2026}

\begin{document}
\maketitle
\begin{abstract}
Structural Entropy Dark Energy (SEDE) reproduces the cosmological data
with no fitted dark-sector parameter, at the price of a single
foundational input: the cosmic apparent horizon carries the
coarse-grained entanglement entropy of its \textbf{volume}, S ∝ V ∝ R³
(the Barrow endpoint Δ = 1), rather than the Bekenstein--Hawking area
law. We ask how much of that input is genuinely irreducible. We show the
postulate decomposes into four sub-claims --- \emph{state, form, scale,
count} --- of which the first three are first-principles or reducible (a
maximal scrambler is volume-law-entangled; volume-law entanglement is
the generic behaviour of a thermal reference state; the magnitude is the
Cohen--Kaplan--Nelson scale), leaving the \textbf{count} (does the
horizon entropy count bulk or boundary degrees of freedom?) as the sole
residue. A Bekenstein no-go proves the count cannot follow from any
energy bound. We then reduce the count to a dynamical statement: the
count is fixed by the \emph{connectivity} of the horizon degrees of
freedom, connectivity by the \emph{range} of the interaction, and
gravity (1/r, strongly long-range) places the horizon in the
non-additive (volume) class --- with area-law arising only for an
isolated, \emph{equilibrium} horizon (a black hole), while the cosmic
horizon is continuously \emph{driven} by structure formation. This
converts the black-hole/cosmic asymmetry into a discriminator and yields
a \textbf{driven non-equilibrium-steady-state (NESS) conjecture}: a
strongly-long-range, structure-driven horizon self-organises to its
area↔volume spinodal and locks the volume branch. We supply both
microscopic maps the conjecture requires --- the cooperative coupling is
the gravitational binding itself
\((J = \lambda _{\mathrm{max}}(W_{\mathrm{grav}}) \geq  J_{\mathrm{c}}\),
generic at the horizon's N), and the drive is the structure-deposited
entropy (clearing the barrier at z ≈ z*) --- and verify the volume-lock
is robust over a broad, untuned basin. Finally we make ``closure''
precise: three routes (double-scaled SYK, Euclidean de Sitter saddles, a
thermal free energy) converge on the same irreducible statement, the de
Sitter static-patch state count, with a decisive empirical settler (Δ
via DESI DR3 + Euclid) already in hand. The volume-law postulate is
thereby narrowed from an unconstrained assumption to a single,
sharply-posed, falsifiable statement.
\end{abstract}

{
\hypersetup{linkcolor=}
\setcounter{tocdepth}{2}
\tableofcontents
}
\section{1. Introduction}\label{introduction}

Horizon thermodynamics is the most robust pointer we have toward the
microscopic structure of gravity. The first law on a causal horizon
yields the Einstein equations {[}Jacobson 1995{]}, the apparent-horizon
first law yields the Friedmann equations {[}Cai \& Kim 2005{]}, and the
Bekenstein--Hawking area law S = A/4 organises black-hole and
cosmological thermodynamics alike. Yet the area law is an assumption
about a \emph{state}: it is the entanglement entropy of a particular
(ground-like) horizon state, and generalisations --- Tsallis--Cirto
non-extensive entropy, Barrow's fractal-horizon
\(S \propto  A^{1+\Delta /2}\) {[}Barrow 2020{]}, and the holographic
dark-energy programme built on them {[}Saridakis 2020{]} --- explore
what happens when that assumption is relaxed.

Structural Entropy Dark Energy (SEDE; Paper I) is one such model,
distinguished by sourcing dark energy from horizon entropy \emph{gated
by structure growth}, \(\rho _{\mathrm{DE}} = T_{\mathrm{AH}}\)
\(s_{\mathrm{grav}}\) \(f_{\mathrm{sat}}(z)\). Its appeal is that the
entire dark sector is parameter-free once \(\Omega _{\mathrm{m}}\) and
one input are fixed: the H-coupling λ = 1−Δ/2 = ½, the equation of
state, the structure coupling γ, and the sound speed all follow. The one
input is the \textbf{volume-law postulate}: the horizon entropy is the
volume-law (Δ = 1) entanglement entropy, not the area law. Paper I
states plainly that this is the model's single irreducible assumption,
motivated but not derived, and that it is falsifiable --- Δ is measured
by upcoming surveys.

This paper asks the prior question: \emph{how irreducible is it,
really?} We are not trying to solve quantum gravity; we are trying to
determine the minimal, sharply-posed statement that the cosmology
actually rests on, and to derive everything reducible to it. The result
is a chain that trades a static, untestable counting postulate for a
dynamical, falsifiable self-organisation statement, with a clean
hand-off to an open --- but well-posed --- problem in de Sitter
holography. Several steps borrow established machinery: the
Cohen--Kaplan--Nelson holographic energy bound {[}Cohen, Kaplan \&
Nelson 1999{]} for the scale; the Page curve and eigenstate
thermalisation for the state; the Campa--Dauxois--Ruffo classification
of long-range systems {[}Campa, Dauxois \& Ruffo 2009{]} for the
connectivity; Ryu--Takayanagi min-cut for the count↔connectivity link;
and the double-scaled-SYK/de Sitter correspondence {[}Susskind 2021;
Narovlansky \& Verlinde 2023{]} for the closure attempt. The
contribution is the chain, and the precise localisation of what remains
open.

Throughout, every quantitative claim is a runnable experiment; the
inline \texttt{run\_*.py} names and the ledger in §8 give the
correspondence, and \texttt{reproduce\_all.py} runs them all with
validation assertions.

\begin{figure}
\centering
\pandocbounded{\includegraphics[keepaspectratio]{../../output/foundations_fig1_reduction.pdf}}
\caption{Fig1}
\end{figure}

\textbf{Figure 1.} The reduction at a glance. The volume-law postulate
splits into \emph{state}, \emph{form}, \emph{scale} and \emph{count};
the first three are derived or reduced (green/blue) and the count is the
residue (orange). The count reduces through connectivity and interaction
range to a driven non-equilibrium steady state, ratcheted through the
spinodal at z* and locked at volume, leaving a single dS-holography
question that the Δ measurement decides empirically. Each box names the
experiment that establishes it.

\subsection{1.1 Relation to prior work}\label{relation-to-prior-work}

SEDE sits at the confluence of several programmes, and its closest
relative deserves explicit positioning.

\textbf{Entropy-production cosmic acceleration (GREA).} The General
Relativistic Entropic Acceleration theory of García-Bellido and
Espinosa-Portalés {[}arXiv:2106.16014, 2106.16012{]} is the nearest
existing model: it too abandons Λ, sources the late-time acceleration
from cosmic-\emph{horizon} entropy, takes the Helmholtz free energy F =
U − TS (not matter alone) as the gravitational source, and makes the
breaking of time-reversal by \emph{entropy production} the engine of
acceleration. We regard GREA as prior art for ``acceleration from
horizon entropy production without Λ'', and our driven-NESS picture
(§4--6) is conceptually continuous with it. The differences are sharp
and testable: (i) GREA uses the \textbf{boundary (area)} horizon entropy
and a single parameter α (the horizon-to-curvature ratio), whereas
SEDE's defining claim is the \textbf{volume} entropy (Δ = 1) ---
precisely the counting residue this paper isolates; (ii) GREA's entropy
growth is driven by the horizon's own \emph{expansion}, whereas SEDE's
is \emph{gated by structure formation} (the \(f_{\mathrm{sat}}\)
factor), which locks w(z) to the growth history and yields a phantom
crossing rather than GREA's quintessence-like w(a); (iii) SEDE's dark
sector is parameter-free. Notably, were SEDE's residue to resolve to the
\emph{area} branch (Δ = 0), its phenomenology would collapse toward
GREA's --- so the Δ measurement (§7) that decides SEDE's residue also
discriminates the two models. A holographic reading of GREA has recently
appeared {[}arXiv:2511.19546{]}, and a broader open-system view of
gravitational non-equilibrium thermodynamics is developing
{[}e.g.~arXiv:1111.6510, 2507.03103{]}; SEDE's contribution to that line
is the reduction of the volume-law \emph{counting} input to a driven
steady state. The connection is also quantitative: GREA's entropy
production maps to an effective bulk viscosity with negative pressure,
and SEDE's structure-driven entropy production d ln
\(f_{\mathrm{sat}}/d\) ln a \textgreater{} 0 maps to a \emph{positive}
(second-law-respecting) effective bulk viscosity
\(\zeta (z) = \rho _{\mathrm{DE}}\) (d ln \(f_{\mathrm{sat}}/d\) ln
a)/(9H²) that peaks at the structure-formation epoch rather than
tracking the horizon expansion (\texttt{run\_lit\_inferences.py}); the w
= −1 crossing is the GREA negative-pressure balance, structure-timed.
SEDE is, in this sense, a \emph{structure-gated} member of the
GREA/entropic-dark-energy family --- the gate being exactly what makes
its w(z) lock to the growth history (the §5 phase-lock) rather than to
the horizon's own size.

\begin{figure}
\centering
\pandocbounded{\includegraphics[keepaspectratio]{../../output/foundations_fig5_inferences.pdf}}
\caption{Fig5}
\end{figure}

\textbf{Figure 2.} Two consequences drawn from the new literature.
\emph{(a)} The de Sitter horizon has negative specific heat (C = −2S);
with gravity strongly long-range this makes the ensembles inequivalent
and the canonical (area-law) accounting ill-defined, selecting the
microcanonical (volume) branch (§4). \emph{(b)} SEDE's structure-driven
entropy production maps to a positive (second-law) effective bulk
viscosity ζ(z) that peaks at the gate-activation epoch (z ≲ 0.5; it
tracks d ln \(f_{\mathrm{sat}}/d\) ln a, not the z≈2 star-formation
peak), casting SEDE as a structure-gated member of the
GREA/entropic-dark-energy family; w(z) crosses −1 at the
viscous--dilution balance.

\textbf{Holographic dark energy and generalised entropy.} SEDE is, at
the background level, a structure-gated Barrow holographic dark energy:
the CKN bound {[}Cohen, Kaplan \& Nelson 1999{]} fixes the scale,
Barrow's fractal-horizon entropy {[}Barrow 2020{]} and its
holographic-DE use {[}Saridakis 2020{]} fix the form, and the
Tsallis--Cirto non-extensive entropy {[}Tsallis \& Cirto 2013{]}
supplies the δ ↔ Δ map. Our addition is the \emph{derivation} question
--- which of these inputs is irreducible --- and the answer that only
the count is.

\textbf{Horizon thermodynamics and entanglement.} The first-law
derivations of Einstein {[}Jacobson 1995{]} and Friedmann {[}Cai \& Kim
2005{]} equations underlie the construction; the
entanglement-equilibrium refinement {[}Jacobson 2016{]} is the basis of
route B; the Ryu--Takayanagi min-cut {[}Ryu \& Takayanagi 2006{]}
supplies the count ⟺ connectivity link; and the volume-law = thermal
identification rests on eigenstate thermalisation {[}Deutsch 1991;
Srednicki 1994; Rigol et al.~2008{]}.

\textbf{de Sitter holography.} The residue --- ``dim 𝓗(static
\(patch) = e^{Area}\) or \(e^{Vol}\)?'' --- is posed in the language of
finite-dimensional de Sitter holography {[}Banks 2001{]} and is the
target of the double-scaled-SYK/de Sitter correspondence {[}Susskind
2021; Narovlansky \& Verlinde 2023; Lin 2022{]}, which we use in the
closure attempt (§7). The maximal-scrambler state invokes the chaos
bound {[}Maldacena, Shenker \& Stanford 2015{]} and SYK random-matrix
universality {[}Cotler et al.~2017{]}.

\textbf{Many-body and complexity inputs.} The long-range/non-additivity
classification is {[}Campa, Dauxois \& Ruffo 2009{]}; the
self-organisation framing is Bak--Tang--Wiesenfeld self-organised
criticality {[}Bak, Tang \& Wiesenfeld 1987{]}, whose cosmological
applications to date concern the \emph{primordial} universe and the
dark-matter sector rather than the late-time horizon. The prototype's
free-fermion \emph{negative} (volume-law needs interactions, not free
transport) is consistent with the monitored free-fermion entanglement
literature {[}Fisher et al.~2023{]}, and that genuine \textbf{volume-law
non-equilibrium steady states} exist is established in driven many-body
systems {[}Ippoliti, Rakovszky \& Khemani 2022{]} --- a condensed-matter
precedent for, not a derivation of, the driven-NESS we invoke.

\textbf{Structure-sourced dark energy.} The general idea that dark
energy is sourced by cosmic structure predates SEDE in several distinct
mechanisms: Gough's information dark energy (≈2008); the
spacetime-fragmentation/``FRW-islands'' proposal of Hossain (2007); and
the cosmological-backreaction and averaging programme (Buchert;
Wiltshire's timescape). SEDE Paper I credits and distinguishes these.
What is specific to SEDE --- and, to our knowledge, without precedent
--- is the \emph{combination} the present paper analyses: the source is
the \textbf{volume-law horizon entropy} (Δ = 1, a counting claim), it is
\textbf{driven/self-organised by structure} through the
\(f_{\mathrm{sat}}\) gate, and the dark sector is
\textbf{parameter-free}. The closest entropic relative remains GREA
(above), which differs in using the boundary (area) entropy. A
systematic priority search returned no work combining these three
elements; this companion concerns that horizon-entropy foundation.

\section{2. The postulate, decomposed}\label{the-postulate-decomposed}

``Volume-law'' conflates four claims of very different status:

\begin{itemize}
\tightlist
\item
  \textbf{state} --- the horizon degrees of freedom are maximally
  entangled (thermal), not in a low-entanglement ground state;
\item
  \textbf{form} --- the entropy scales as the volume, S ∝ V;
\item
  \textbf{scale} --- the magnitude is \(\rho _{\mathrm{DE}}\) ∼
  \(\rho _{\mathrm{crit}}\), not \(\rho _{\mathrm{Planck}}\);
\item
  \textbf{count} --- the \emph{number} of horizon dof grows as the bulk
  \((N \propto  R^{d-1})\) rather than the boundary
  \((N \propto  R^{d-2})\).
\end{itemize}

Since S ∼ (state factor) × N and the Barrow deformation enters as
\(S \propto  A^{1+\Delta /2}\)∝ \(R^{(d-2)(1+\Delta /2)}\), the exponent
Δ is fixed by the \textbf{count} (area → Δ=0, volume → Δ=1 in d=4); the
state only sets whether the prefactor is maximal. So the postulate's
empirical content --- Δ --- is the counting claim, and the other three
are separable. The rest of the paper derives or reduces state, form and
scale, and isolates count.

\section{3. How derivable is the postulate? The route
map}\label{how-derivable-is-the-postulate-the-route-map}

Five routes (\texttt{run\_deriv\_\{A..E\}\_*.py}):

\textbf{A --- de Sitter holography / maximal scrambler {[}DERIVED state;
OPEN count{]}.} A Majorana-SYK diagonalisation confirms the horizon
\emph{state} is maximally entangled / volume-law: chaotic level
statistics (⟨r⟩ ≈ 0.71) and Page-value saturation
\((S/S_{\mathrm{Page}} \approx  0.99)\). But SYK is all-to-all /
geometry-free and cannot pose the bulk-vs-boundary count; the count is
the genuine open dS-holography sub-target.

\textbf{B --- entanglement first law in a thermal state {[}REDUCED
form{]}.} Jacobson's derivation gets the area-law Einstein equations
from maximal entanglement of the \emph{vacuum}. Redone around a
\emph{thermal} (de Sitter Gibbs) reference, the entanglement entropy
acquires a volume term that dominates the area term once the region
exceeds the thermal length \(\lambda _{\mathrm{th}} = 1/T\). Volume-law
is thus the generic thermalised behaviour, not exotic. At the de Sitter
temperature alone the horizon sits just on the area side
\((R_{\mathrm{H}}/\lambda _{\mathrm{th}} = 1/2\pi )\); volume-law
dominance requires thermalisation above the de Sitter bath (toward
Planck = maximal scrambling), whereupon the density is Planckian --- the
CKN scale. So the \emph{form} reduces to thermalisation, leaving the
\emph{scale} to CKN.

\textbf{C --- Verlinde emergent gravity {[}REDUCED, unification{]}.} A
volume-law de Sitter entanglement entropy is exactly what Verlinde's
emergent-gravity programme invokes for apparent dark matter; one
volume-law scale gives both \(\rho _{\mathrm{DE}}\) ∼
\(\rho _{\mathrm{crit}}\) and a₀ = cH₀/2π (≈ 0.87× the
radial-acceleration-relation value). A unification motivation that ties
the postulate to an independent anomaly; it inherits that programme's
open issues.

\textbf{D --- gravitational non-additivity {[}NEGATIVE{]}.} The
Tsallis--Cirto map δ = 1 + Δ/2 makes Δ=1 ⟺ δ=3/2 exact, but pinning δ
from a measurable non-additivity fails: real cluster kinematics are
Gaussian (q ≈ 1.01, \texttt{run\_cluster\_tsallis.py}), not q ≈ 1.5 ---
the weakest route.

\textbf{E --- no-go from energy bounds {[}DERIVED{]}.} The Bekenstein
bound on the horizon energy gives S ≤ 2πRE/ħc =
¼\((A/\ell _{\mathrm{P}}^{2})\) --- \emph{exactly} the area law
(verified \(S_{\mathrm{Bek}}/S_{\mathrm{area}} = 1.000\), the known
10¹²²). The volume law therefore cannot come from any
maximum-entropy/energy argument; it exceeds the bound by precisely
\(R/\ell _{\mathrm{P}}\) ∼ 10⁶¹ (the ``seam''). This proves the missing
ingredient is a \textbf{state/counting} input and fixes its size.

\emph{Net:} state (A) and form (B) and scale (CKN) are first-principles
or reduced; the irreducible residue is the bulk-vs-boundary
\textbf{count}.

\section{4. Reducing the residue to a driven
NESS}\label{reducing-the-residue-to-a-driven-ness}

The count is not a free coin-flip --- but here we must be careful to
avoid a circularity, because the entanglement class cannot simply be
\emph{posited}. It is fixed by the Hamiltonian and the state, and the
decisive fact is that \textbf{the entanglement area law is a theorem
with a hypothesis}: it is proved for \emph{local} (finite-range)
Hamiltonians --- Hastings (2007) rigorously in 1D via Lieb--Robinson
bounds, Brandão--Horodecki (2015) from exponential decay of
correlations, reviewed in Eisert--Cramer--Plenio (2010) --- and the
proof \emph{requires} exponential clustering. Gravity is 1/r (α = 1),
strongly long-range (α ≤ d) and non-additive {[}Campa, Dauxois \& Ruffo
2009{]}, and there the hypothesis fails: in an explicit gapped
free-fermion model we find the correlation function decays as a
\emph{power law} for α = 1 (slope −1.5) versus \emph{exponentially} for
the short-range reference (slope −6.9), exactly as Koffel--Lewenstein--
Tagliacozzo (2012) found (\texttt{run\_v1\_locality.py}). So the
area-law theorem \textbf{does not apply} to a gravitationally-bound
horizon: area-law is \emph{not guaranteed}, and the count is genuinely
open between area and volume with neither the default. We are explicit
about what this does and does not give. It does \emph{not} by itself
derive volume --- a noncritical long-range ground state can still be
area-law {[}Kuwahara--Saito 2020{]}, and indeed the gapped ground-state
block entropy stays bounded for every α we test. The volume law
additionally requires the horizon \emph{state} to be thermal / maximally
scrambled (route A), where entanglement is volume-law; and even that
fixes the \emph{state}, not the \emph{count}. The super-extensive
non-additivity (\texttt{run\_residue\_longrange.py};
\(u \propto  N^{1-\alpha /d}\), slopes 0.53 in d=2, 0.68 in d=3 vs 0.50,
0.67) is the thermodynamic face of the same long-range property. The
honest conclusion is that volume is the \emph{motivated, non-default}
class for a self-gravitating thermal horizon --- area-law arises only
when the holographic bound intervenes for an \emph{isolated,
equilibrium} horizon (a black hole saturates the Bekenstein bound, route
E, and relaxes to area-law) --- while the bulk-vs-boundary
\textbf{count} itself remains the open dS-holography residue. One caveat
keeps this honest: the area-law theorems are statements about many-body
systems with a specified Hilbert space and Hamiltonian, and what the
horizon's degrees of freedom \emph{are} --- and what Hamiltonian governs
them --- is precisely the open dS-holography question. So this is a
motivated \emph{conditional} --- \emph{if} the horizon is a long-range
many-body system, \emph{then} area-law is not forced and volume is
favoured --- not a theorem about the actual horizon. What it establishes
is that the entanglement class is read off the (modelled) Hamiltonian
rather than assumed --- which is what defeats the circularity charge ---
not that the count is derived.

This same long-range property has a second, independent consequence that
bears on the canonical-vs-microcanonical λ ambiguity (whether
\(\rho _{\mathrm{DE}} \propto  H^{2}\) \(f_{\mathrm{sat}}\), area, or ∝
H \(f_{\mathrm{sat}}\), volume). A non-additive system has
\emph{inequivalent} statistical ensembles {[}Campa, Dauxois \& Ruffo
2009{]}, the signature being negative specific heat --- and the de
Sitter horizon indeed has C = −2S \textless{} 0 throughout its history
(\texttt{run\_lit\_inferences.py}). For such a system the canonical
ensemble is ill-defined; the isolated self-gravitating horizon must be
treated \emph{microcanonically} (by its state count). The canonical
accounting is precisely the area-law (λ = 1) branch, so ensemble
inequivalence removes it as unphysical and selects the microcanonical
(volume, λ = 1/2) branch the model uses. This does not close the residue
--- whether the microcanonical state count is itself volume is the
dS-holography question --- but it eliminates one of its two horns from
first principles. The cosmic horizon differs in one respect: it is
continuously \emph{driven} by structure formation (the
\(f_{\mathrm{sat}}\) gate, \(dS_{\mathrm{struct}}/dt > 0\); present
drive ≈ 0.45). This turns the black-hole/cosmic asymmetry into a
\textbf{discriminator} --- equilibrium → area, driven NESS → volume ---
and reduces the residue to a single, falsifiable statement:

\begin{quote}
\textbf{Driven-NESS conjecture.} A strongly-long-range horizon, driven
by structure formation, settles in the volume-law steady state rather
than relaxing to the area-law equilibrium.
\end{quote}

\section{5. The driven-NESS conjecture and its two
maps}\label{the-driven-ness-conjecture-and-its-two-maps}

\textbf{Prototype (\texttt{run\_driven\_ness.py}).} Generic
driven-dissipative free fermions do \emph{not} realise the effect
(uniform drive is trivially extensive and connectivity-independent;
boundary drive decays to area-law) --- a NEGATIVE that sharpens the
content to \emph{cooperative bistable locking}. With the area↔volume
order parameter under cooperative (Wilson--Cowan) dynamics, a transient
structure-drive pulse locks the volume branch (m → 0.93, persistent) for
long-range cooperativity \((J > J_{\mathrm{c}})\), while short-range
tracks-then-relaxes to area and the undriven case stays at the area
ground. Only drive × long-range selects volume.

\textbf{Map 1 --- connectivity → \(J \geq  J_{\mathrm{c}}\)
{[}DERIVED{]} (\texttt{run\_gravitational\_coupling.py}).} The
cooperative coupling is the largest eigenvalue of the gravitational
coupling matrix,
\(J_{\mathrm{eff}} = \lambda _{\mathrm{max}}(W_{\mathrm{grav}})\),
\(W_{\mathrm{ij}} = g/r_{\mathrm{ij}}^{\alpha }\). For a homogeneous
horizon this equals the per-site gravitational binding
\((\lambda _{\mathrm{max}}/row\)-sum = 1.02) --- the \emph{same}
super-extensive quantity behind the non-additivity of §4. So
volume-counting and volume-locking are one number. It scales as
\(N^{1-\alpha /d}\) (slopes 0.52, 0.67; a short-range α\textgreater d
control is bounded, \(N^{0.01}\)), so since α=1≤d it exceeds any fixed
\(J_{\mathrm{c}}=4T\) at the horizon's N ∼ 10¹²² --- bistability is
generic, not tuned.

\begin{figure}
\centering
\pandocbounded{\includegraphics[keepaspectratio]{../../output/criticality_soc.pdf}}
\caption{Fig2}
\end{figure}

\textbf{Figure 3.} Why the criticality is the volume law, not a separate
assumption. \emph{(a)} The area↔volume free energy is a smooth crossover
for the bare gate (J = 0) and develops a spinodal only for
\(J \geq  J_{\mathrm{c}}\). \emph{(b)} The bare-gate susceptibility is
finite; the spinodal diverges at the operating point m = ½. \emph{(c)}
The cooperative coupling
\(J = \lambda _{\mathrm{max}}(W_{\mathrm{grav}})\) is the gravitational
binding itself: all-to-all (volume-law) connectivity gives
\(\chi _{\mathrm{peak}} \propto  \surd N\) → ∞ (a real spinodal), while
local (area-law) connectivity is bounded.

\textbf{Map 2 --- deposition → drive {[}DERIVED{]}
(\texttt{run\_deposition\_drive.py}).} The drive is the accumulated
structure-deposited horizon entropy (the collapsed
\(fraction \propto  f_{\mathrm{sat}}\)). Its timing is derived --- the
drive clears the area-well barrier \(h_{\mathrm{c}}\) when
\(f_{\mathrm{sat}} \approx  h_{\mathrm{c}}\), i.e.~at z ≈ 1.45 ≈ z* ---
and integrating the order-parameter dynamics with this physical drive
flips the horizon area→volume near z* and locks it (Δ = 1 permanent),
whereas an un-amplified drive never leaves area.

\begin{figure}
\centering
\pandocbounded{\includegraphics[keepaspectratio]{../../output/deposition_drive.pdf}}
\caption{Fig3}
\end{figure}

\textbf{Figure 4.} The deposition → drive map. The drive is the
accumulated structure entropy (∝ \(f_{\mathrm{sat}}\)); it crosses the
area-well barrier \(h_{\mathrm{c}}\) at z ≈ z*, flipping the order
parameter to volume and locking it there permanently (red), whereas the
un-amplified bare deposition stays on the area branch (blue).

\textbf{The drive is a violent-relaxation deposit, and the γ-mechanism
is the same event} (\texttt{run\_violent\_relaxation.py}). Structure
forms by gravitational collapse, i.e.~by \emph{violent relaxation}
{[}Lynden-Bell 1967{]}, which drives a long-range system on the
dynamical time to a quasi-stationary state rather than to thermal
equilibrium. This single event does two things SEDE otherwise treats
separately. First, it dissipates the halo binding energy
\(E_{\mathrm{bind}} = M\sigma _{\mathrm{v}}^{2} \propto  M^{5/3}\)
(virial: \(R_{\mathrm{vir}} \propto  M^{1/3}\),
\(\sigma _{\mathrm{v}}^{2} \propto  M^{2/3}\)) into the horizon ledger
--- the mass weight that \emph{is} the structure coupling γ ≈ 1.5 (Paper
I, Theorem 4C). Second, that same deposited binding energy (∝
⟨\(\sigma _{\mathrm{v}}^{2}\)⟩) is the amplitude of the NESS drive
above. So the γ-weight and the driven-NESS drive are one quantity from
one process --- the volume lock is the horizon's Lynden-Bell
quasi-stationary state, fed by violent relaxation. Two corollaries
follow at no cost: collapse is prompt
\((t_{\mathrm{dyn}}/t_{\mathrm{Hubble}}\)=
\(\Delta _{\mathrm{vir}}^{-1/2} \approx  0.07\)), so the horizon is
driven by an ongoing sequence of such events whose envelope is
\(f_{\mathrm{sat}}\); and the quasi-stationary state's lifetime in a
long-range system scales as \(N^{\delta }\) {[}Campa, Dauxois \& Ruffo
2009{]}, so at the horizon's N ∼ 10¹²² the lock is permanent --- a
second permanence argument beside the second-law one (route E). Both the
halo QSS and the de Sitter horizon have negative specific heat,
consistent with the microcanonical (volume) selection of §4.

\section{\texorpdfstring{6. Driven \emph{through} the spinodal: ratchet
and
robustness}{6. Driven through the spinodal: ratchet and robustness}}\label{driven-through-the-spinodal-ratchet-and-robustness}

Developing the conjecture corrected its statement
(\texttt{run\_soc\_attractor.py}): the system locks at the \emph{volume}
well, not \emph{at} the spinodal. Because the structure
\(drive \propto  f_{\mathrm{sat}}\) is monotone, it \emph{necessarily}
climbs to the field at which the area well vanishes --- the area-branch
spinodal --- at z ≈ z\emph{; the crossing is a forced \textbf{ratchet},
and the generalised second law then locks the volume well (volume is the
global entropy maximum, \(S_{\mathrm{vol}} \gg  S_{\mathrm{area}}\) by
\(R/\ell _{\mathrm{P}}\), route E). The near-critical fluctuation
signatures are therefore a }transient* at z\emph{, not a permanent
state. The self-organisation content is a \textbf{robustness} claim,
verified: volume-locking holds over a broad contiguous basin of the
landscape parameters \((J \geq  J_{\mathrm{c}}\), θ), over drive
amplitudes 0.6--3×, and from all sub-barrier initial conditions --- an
attractor, not a knife-edge (it fails only at strong holographic pull).
What is }not* derived is the absolute horizon free energy F(m); only
that a broad, untuned range works. So the residue's sharpest form is
``the horizon area↔volume free energy is bistable with
\(J \geq  J_{\mathrm{c}}\) (gravity supplies this) and a sub-unity
barrier'' --- strictly weaker, more physical, and falsifiable than
``assume volume counting''.

\textbf{A new, sharp prediction: the crossing carries mean-field SOC
statistics} (\texttt{run\_soc\_signature.py}). Because the system is
driven \emph{through} its spinodal at z* (not parked at a tuned critical
point) and the coupling is all-to-all/mean-field (§5), the fluctuations
at the crossing are scale-free in the \textbf{mean-field
self-organised-criticality} universality class {[}Bak, Tang \&
Wiesenfeld 1987{]}. That class is parameter-free: the avalanche-size
distribution is \(P(s) \propto  s^{-3/2}\) (we recover τ = 1.51 from the
critical branching process), and the susceptibility and autocorrelation
time diverge as the order parameter crosses the spinodal --- a
\emph{transient} critical slowing-down localised at z* (we find the peak
at z ≈ 1.08 ≈ z\emph{), reddening the temporal spectrum toward 1/f.~This
turns the companion-paper CHR layer (Paper I §6, predictions P3--P5)
from ``an enhancement at z}'' into a specific, falsifiable statement:
any z\emph{-localised excess in growth/ISW should be \textbf{scale-free
with exponent τ ≈ 3/2 and a 1/f spectrum}, generic (self-organised, not
tuned), and }transient\emph{. This discriminates SEDE from clustering
dark energy / k-essence (which carries a single characteristic scale,
the dark-energy sound horizon, not a scale-free spectrum) and from ΛCDM
(no z}-localised feature), and is within reach of redshift-resolved
weak-lensing tomography and ISW cross-correlation (Euclid/Rubin). We
flag the status honestly: this is a \emph{conditional} prediction. The
exponent τ ≈ 3/2 is parameter-free, but the \emph{amplitude} of the
z\emph{-localised excess is not yet forecast --- it is set by the
deposited-entropy fraction \(\epsilon _{\mathrm{dep}}\) ∼ 10⁻⁷ and the
near-critical susceptibility, and a quantitative amplitude and survey
signal-to-noise are deferred to the companion CHR analysis. So τ ≈ 3/2
is a falsifiable shape }conditional on the CHR layer being detectable at
all*, not a claim that it is.

\begin{figure}
\centering
\pandocbounded{\includegraphics[keepaspectratio]{../../output/foundations_fig4_socsignature.pdf}}
\caption{Fig4}
\end{figure}

\textbf{Figure 5.} The new prediction. \emph{(a)} The near-spinodal
fluctuations fall in the mean-field self-organised-criticality class,
whose avalanche-size law is the parameter-free
\(P(s) \propto  s^{-3/2}\) (recovered as τ = 1.51 from critical
branching). \emph{(b)} As the structure drive carries the order
parameter to the spinodal, the area-phase susceptibility and
autocorrelation time diverge --- a \emph{transient} critical
slowing-down localised at z\emph{, reddening the spectrum toward
1/f.~Any z}-localised growth/ISW excess should therefore be scale-free
with τ ≈ 3/2, which single-scale clustering dark energy cannot mimic.

\section{7. What would make it a
closure}\label{what-would-make-it-a-closure}

A closure derives F(m) itself, reducing the dark sector to
\(\Omega _{\mathrm{m}} + established\) physics. The three routes
(\texttt{run\_closure\_attempts.py}) converge:

\begin{itemize}
\tightlist
\item
  \textbf{A (DSSYK).} Entropy extensive in the dof count
  \((S_{\mathrm{max}} = (N/2)ln\) 2 ∝ N, slope 1.00); cannot be
  super-extensive \((volume \propto  N^{3/2})\). We state this plainly:
  \textbf{read in the standard Narovlansky--Verlinde dictionary, the
  leading concrete model of de Sitter holography maps to the de Sitter
  \emph{area} (Gibbons--Hawking) --- i.e.~it currently favours the area
  branch, a real tension SEDE's volume postulate must eventually
  overcome.} The mitigating point --- that DSSYK is all-to-all /
  geometry-free and so cannot, by itself, \emph{pose} the
  bulk-vs-boundary counting question --- is secondary, not a
  neutraliser. As things stand, route A is evidence to be answered, not
  used.
\item
  \textbf{B (Euclidean dS saddles).} Einstein and Barrow horizon free
  energies each have a single saddle --- no volume-law saddle of the
  bare path integral (route-E no-go in saddle form).
\item
  \textbf{C (thermal F(m)).} The bistable landscape \emph{is} derivable
  --- F(m)/T = −(J/2)m² + θm + {[}m ln m + (1−m)ln(1−m){]} from
  gravitational cooperativity
  \((J = \lambda _{\mathrm{max}} \geq  J_{\mathrm{c}})\), the
  holographic pull (route E), and configurational entropy: area well m ≈
  0.07, volume well m ≈ 0.93 whenever \(J \geq  J_{\mathrm{c}}\).
  Level-2 existence, conditional on the volume branch being accessible.
\end{itemize}

All three reduce to one statement --- \emph{the horizon's dof count is
the bulk (volume) one} --- which A and B cannot supply and C uses.
Closure has three precise forms: \textbf{Level 1}, a dS-holography
computation of the bulk count (open, not shortcut-able past route E);
\textbf{Level 2}, exhibiting the volume branch as a real saddle/state (C
gives the landscape conditional on it); \textbf{Level 3},
\emph{measuring} Δ (DESI DR3 + Euclid to ∼0.09; Δ=1 establishes, Δ=0
refutes --- decisive, in hand).

\subsection{7.1 The count is state-dependent --- and that resolves the
``everything points to area''
tension}\label{the-count-is-state-dependent-and-that-resolves-the-everything-points-to-area-tension}

Routes A and B, and the standard handles generally (the de Sitter
modular Hamiltonian is a \emph{local} boost; DSSYK in the
Narovlansky--Verlinde dictionary; the CKN holographic bound; black-hole
thermodynamics), all favour the \textbf{area} count. Read as universal
statements about ``the'' horizon count they look like a standing
objection. The resolution is that they are not universal: \textbf{the
count is state-dependent}, and every one of those handles is an
\emph{equilibrium} horizon. The equilibrium-vs-driven discriminator of
§4 is exactly this --- the \(f_{\mathrm{sat}}\) gate \emph{is} the
fraction of the volume capacity that structure-driving activates,
\(N_{\mathrm{eff}} = N_{\mathrm{area}} + (N_{\mathrm{vol}} - N_{\mathrm{area}})\cdot f_{\mathrm{sat}}\),
so the equilibrium limit \((f_{\mathrm{sat}}\) → 0) is area and the
driven, locked limit \((f_{\mathrm{sat}}\) → 1) is volume
(\texttt{run\_residual\_resolution.py}). (This is \emph{not} the
rejected running-Δ background of §3: the cosmic horizon is driven past
the spinodal early and \emph{locks} at constant Δ = 1; the
state-dependence is \emph{across horizons}, driven cosmic vs equilibrium
black hole, not a running Δ over cosmic time.) So the area results are
SEDE's equilibrium baseline, not counter-evidence. One case needs care
because it is \emph{not} obviously equilibrium: the DSSYK correspondence
is a duality with the de Sitter static patch --- the cosmological
horizon itself. The point is that it is a duality with \emph{eternal}
(matter-free) de Sitter, an equilibrium spacetime with no structure to
drive it \((f_{\mathrm{sat}} \equiv  0)\); the real universe's apparent
horizon is embedded in a matter-plus-structure cosmology and is driven.
So ``DSSYK = area'' is the area of \emph{equilibrium} de Sitter ---
again the \(f_{\mathrm{sat}}\) → 0 limit --- and is consistent with a
volume-law \emph{driven} horizon, not a contradiction of it. We are
explicit that the state-dependent count is itself an effective-model
statement: writing \(N_{\mathrm{eff}}\) as a function of
\(f_{\mathrm{sat}}\), and reading a black hole as the
\(f_{\mathrm{sat}} = 0\) case, extends \(f_{\mathrm{sat}}\) from its
cosmological definition (a function of the growth D(z)) to a putative
property of any horizon's driving --- physically motivated, but the same
conditional as V1, one layer out.

This makes the residue \textbf{doubly falsifiable, on two independent
horizons}. (i) The cosmic horizon (driven) should have Δ = 1 --- DESI
DR3 + Euclid. (ii) The black hole (equilibrium) should have Δ = 0 ---
and the observable is primordial-black-hole evaporation: SEDE
\emph{predicts that PBHs evaporate normally} (standard Hawking), because
black holes are undriven, whereas a \emph{universal}-volume theory
forbids it \((T_{\mathrm{B}}/T_{\mathrm{H}}\) ∼ 7×10⁻²¹, evaporation
time × 10⁸¹). The PBH channel therefore tests the \emph{mechanism} ---
it distinguishes driven-NESS SEDE \((\Delta _{\mathrm{BH}} = 0)\) from
universal-volume quantum gravity \((\Delta _{\mathrm{BH}} = 1)\) --- and
the black hole being area-law, far from a problem, is a SEDE prediction.
This is a \emph{future} discriminator, not present-day evidence: SEDE's
\(\Delta _{\mathrm{BH}} = 0\) means standard Hawking evaporation,
consistent with all current data (including \textasciitilde10¹⁵ g PBH
abundance bounds) precisely because it predicts no new black-hole
effect; the channel discriminates only on a positive PBH-evaporation
detection or a robust population statement.

This also yields a new statement about black holes in general:
\textbf{area-law is the \emph{equilibrium} default, not a fundamental
geometric law} --- a black hole is area-law \emph{because it is
undriven} \((f_{\mathrm{sat}}\) → 0), and the cosmic horizon's
uniqueness is a calculation, not an assumption. The most-driven black
hole conceivable is a primordial one \emph{at formation}, where the
entire horizon is built by a collapsing overdensity (violent
relaxation); yet even there the structure entropy deposited is
negligible against the black-hole entropy \emph{created},
\(S_{\mathrm{rad}}/S_{\mathrm{BH}}\) ∼ \((M/M_{\mathrm{P}})^{-1/2}\), so
the formation drive \(f_{\mathrm{form}}\) ∼
\((M_{\mathrm{P}}/M)^{1/2} \ll  h_{\mathrm{c}}\) even after the
near-critical amplification (\texttt{run\_pbh\_driving.py},
\texttt{run\_bh\_new\_physics.py}). Accreting and merging black holes
are weaker still. So SEDE robustly predicts \emph{all} black holes are
area-law \((\Delta _{\mathrm{BH}} = 0)\), and the cosmic apparent
horizon --- uniquely driven \emph{sustainedly} to \(f_{\mathrm{sat}}\) →
1 over the whole structure-formation history --- is the only volume-law
horizon. The falsifier is correspondingly sharp: an \emph{isolated}
volume-law black hole (anomalously slow PBH evaporation, or a
modified-area GW ringdown) would mean Δ is \emph{universal}
(Barrow-type) and would refute the state-dependent,
driven-vs-equilibrium mechanism --- so SEDE is \emph{more} falsifiable
on black holes than a universal-Barrow theory, not less.

The honest theoretical residue is correspondingly narrowed: not ``derive
the count,'' but ``prove that a \emph{driven} horizon activates the
volume capacity'' --- a statement now \emph{consistent with} the
equilibrium area results (their \(f_{\mathrm{sat}}\) → 0 limit), with
the two-channel measurement deciding it either way regardless of dS
holography.

Two final caveats, stated openly. First, the \emph{magnitude}: the
structure drive is \(\epsilon _{\mathrm{dep}}\) ∼ 10⁻⁷, yet the count it
flips spans area to volume --- a factor ∼ \(R/\ell _{\mathrm{P}}\) ∼
10⁶¹ in activated degrees of freedom. The near-critical bistability is
therefore doing heavy lifting: the drive does not \emph{add} the dof, it
\emph{nucleates} the area→volume transition past the spinodal, and the
second law carries the system to the volume well, where it locks
(§5--6). That a 10⁻⁷ trigger produces a 10⁶¹ change in the count is the
V1/§7.1 conditional stated quantitatively, and it rests entirely on the
(modelled) bistable landscape. Second, the \emph{pre-registered} case: a
measured cosmic Δ significantly below 1 --- say Δ ≈ 0.5 from DESI DR3 +
Euclid --- would be evidence \emph{against} the locked-Δ = 1 picture,
pointing either to incomplete volume activation or to a genuinely
intermediate count, and would \emph{not} be a free parameter to absorb;
the model predicts Δ = 1 for the driven cosmic horizon and is falsified
by a robust intermediate value.

\subsection{\texorpdfstring{7.2 The count as a horizon Hausdorff
dimension, and what remains of it
(\texttt{run\_ds\_count.py})}{7.2 The count as a horizon Hausdorff dimension, and what remains of it (run\_ds\_count.py)}}\label{the-count-as-a-horizon-hausdorff-dimension-and-what-remains-of-it-run_ds_count.py}

The Level-1 question --- ``is dim 𝓗(de Sitter static
\(patch) = e^{Area/4}\) or \(e^{Volume}\)?'' --- has a clean geometric
form: dim 𝓗 = \(e^{Area}\) ⟺ the horizon is a smooth 2-surface
(Hausdorff dimension \(d_{\mathrm{H}} = 2\)), and dim 𝓗 = \(e^{Volume}\)
⟺ it is a space-filling fractal \((d_{\mathrm{H}} = 3)\), so the count
\emph{is} the horizon's Hausdorff dimension,
\(\Delta  = d_{\mathrm{H}} - 2\). The standard frameworks all give a
smooth horizon: the Gibbons--Hawking maximum, the finite-dimensional
proposal {[}Banks 2001{]}, the Type II₁ observer algebra in which
\(S_{\mathrm{dS}}\) is the maximal entropy {[}Chandrasekaran, Longo,
Penington \& Witten 2023{]}, and DSSYK. There is, however, a concrete
mechanism by which the \emph{driven} horizon could differ: a horizon
driven by a stochastic source is a roughening 2-surface, and a
self-affine surface of roughness exponent α has
\(d_{\mathrm{H}} = 3 - \alpha\), hence Δ = 1 − α. The identification we
make is that the roughening \emph{noise is the structure drive} --- so
an undriven horizon is the smooth minimum \((d_{\mathrm{H}} = 2\), area)
and a structure-driven horizon roughens \((d_{\mathrm{H}}\) → 3,
volume). This is the microscopic, horizon-geometry form of the
state-dependent count of §7.1, and it \emph{reconciles} the area
frameworks (which describe the eternal, equilibrium horizon) with SEDE's
volume (the driven horizon) as two limits of one statement, rather than
as a tension. It reduces the residue to two sharp, well-posed questions.
(i) Is the driven roughening \emph{linear} (Edwards--Wilkinson, the 2+1D
lower-critical class, α → 0 ⇒ Δ = 1) or \emph{nonlinear} (KPZ, α ≈ 0.39
⇒ Δ ≈ 0.6)? This is settled by two textbook inputs rather than a
conjecture (\texttt{run\_kpz\_membrane.py}): the KPZ coefficient
\emph{is} the normal growth velocity (λ = v, from the
height-representation geometry ∂h/∂t = v√(1+(∇h)²) ≈ v + (v/2)(∇h)²),
and a driven interface's velocity is proportional to its driving free
energy (v = M·F/V, interface kinetics --- for a horizon, the
membrane-paradigm statement that it grows in response to the free-energy
flux across it {[}Damour 1978; Thorne, Price \& Macdonald 1986{]}). So λ
= M·F/V ∝ F; and SEDE's free energy vanishes identically
\((F = E - T_{\mathrm{AH}}\) S = 0, since
\(\rho _{\mathrm{DE}} = T_{\mathrm{AH}}\) \(s_{\mathrm{grav}}\); §8),
giving λ = 0 ⇒ EW ⇒ Δ = 1. The one physical assumption is that the
horizon grows by membrane-paradigm interface kinetics --- the standard
effective description of horizon dynamics. (ii) Does the marginal 2+1D
roughening count as a genuine \(d_{\mathrm{H}} = 3\) fractal? Here we
must be candid (\texttt{run\_ds\_marginal.py}): the horizon is a
2-surface, which is \emph{exactly} the lower critical dimension of EW
roughening (α = (2−d)/2 = 0 at d = 2) --- below it the same dynamics is
a genuine fractal, above it is smooth, and at d = 2 it is
\emph{marginal}, space-filling only logarithmically. So the roughening
picture does \textbf{not robustly force} Δ = 1; it explains why the
count is delicate --- the horizon sits precisely on the area↔volume
boundary --- and leaves the tie to sub-leading physics. This tempers
``EW ⇒ Δ = 1'' to ``EW ⇒ Δ = 1 \emph{marginally},'' and it means a
measured Δ ∈ (0, 1) (the log-suppressed case) would be \emph{within} the
picture, not a refutation.

The honest bottom line, then, is that we have \textbf{not resolved de
Sitter holography} --- no one has computed dim 𝓗(static patch)
non-perturbatively, and we do not claim to
(\texttt{run\_ds\_resolve.py}). What we have done is turn the count from
an unstructured ``area or volume?'' into a \emph{geometric,
empirically-decidable} statement \((\Delta  = d_{\mathrm{H}} - 2)\),
reduce its theoretical residue to two named obstructions --- the
marginal dimension (O1) and the Planck-scale validity of the effective
descriptions (O2) --- and identify the single open computation an actual
resolution requires. The equilibrium value is Δ = 0 and the driven value
is Δ → 1, marginally; the empirical Δ (DESI DR3 + Euclid, §6) is the
arbiter.

\section{8. Discussion}\label{discussion}

The chain reduces SEDE's one foundational input from a static,
untestable counting postulate to a single dynamical, falsifiable
statement: \emph{the de Sitter static patch's degrees of freedom count
its bulk, not its boundary.} Everything else --- state, form, scale,
both microscopic maps, and the bistable landscape conditional on the
volume branch --- is derived or reduced. Two honest limits bound the
claim. First, this is a \emph{reduction}, not a closure: the volume
branch is used, not derived, and deriving it is the open dS-holography
problem (route E forbids energy-bound shortcuts). Second, the dynamical
model is effective/mean-field; the connectivity→J and deposition→drive
maps are physical, but F(m)'s absolute parameters are not computed from
the microscopic horizon, only shown to lie in a broad untuned basin.
What makes the reduction worthwhile is that it is \emph{falsifiable
where the postulate was not}: it predicts a transient near-critical
epoch at z* and ties Δ to the structure-driving history, and Δ itself is
pinned to ∼0.09 by DESI DR3 + Euclid. The residue is then not a weakness
in the cosmology but a clean hand-off to a well-posed quantum-gravity
question, with a measurement that settles it either way.

\subsection{Ledger}\label{ledger}

\begin{longtable}[]{@{}
  >{\raggedright\arraybackslash}p{(\linewidth - 6\tabcolsep) * \real{0.2500}}
  >{\raggedright\arraybackslash}p{(\linewidth - 6\tabcolsep) * \real{0.2500}}
  >{\raggedright\arraybackslash}p{(\linewidth - 6\tabcolsep) * \real{0.2500}}
  >{\raggedright\arraybackslash}p{(\linewidth - 6\tabcolsep) * \real{0.2500}}@{}}
\toprule\noalign{}
\begin{minipage}[b]{\linewidth}\raggedright
step
\end{minipage} & \begin{minipage}[b]{\linewidth}\raggedright
result
\end{minipage} & \begin{minipage}[b]{\linewidth}\raggedright
tag
\end{minipage} & \begin{minipage}[b]{\linewidth}\raggedright
script
\end{minipage} \\
\midrule\noalign{}
\endhead
\bottomrule\noalign{}
\endlastfoot
state & maximal entanglement (Page saturation) & DERIVED &
\texttt{run\_deriv\_A\_dsholo.py} \\
form & volume = thermal-state entanglement & REDUCED &
\texttt{run\_deriv\_B\_thermal.py} \\
scale & \(\rho _{\mathrm{DE}}\) ∼ \(\rho _{\mathrm{crit}}\) (CKN) &
DERIVED & --- \\
no-go & energy bounds ⇒ area, not volume & DERIVED &
\texttt{run\_deriv\_E\_nogo.py} \\
residue→range & gravity long-range (α≤d) ⇒ volume class & REDUCED &
\texttt{run\_residue\_longrange.py} \\
conjecture & drive × long-range locks volume & MODELLED &
\texttt{run\_driven\_ness.py} \\
map 1 &
\(J_{\mathrm{eff}} = \lambda _{\mathrm{max}}(W_{\mathrm{grav}}) = binding \geq  J_{\mathrm{c}}\)
& DERIVED & \texttt{run\_gravitational\_coupling.py} \\
map 2 & \(drive \propto  f_{\mathrm{sat}}\) clears barrier at z* &
DERIVED & \texttt{run\_deposition\_drive.py} \\
spinodal & ratchet through spinodal; robust basin & DERIVED (robust) &
\texttt{run\_soc\_attractor.py} \\
closure & three routes ⇒ the dS state count & analysis &
\texttt{run\_closure\_attempts.py} \\
ensemble & C = −2S + long-range ⇒ microcanonical; area horn removed &
REDUCED & \texttt{run\_lit\_inferences.py} \\
unify & γ-mechanism = driven-NESS (one Lynden-Bell QSS) & DERIVED &
\texttt{run\_violent\_relaxation.py} \\
prediction & z*-fluctuations scale-free, τ ≈ 3/2 (mean-field SOC) &
PREDICTION & \texttt{run\_soc\_signature.py} \\
V1 defence & area-law is a \emph{local}-Hamiltonian theorem; gravity
(power-law correlations) voids it & ARGUED &
\texttt{run\_v1\_locality.py} \\
residual & count is state-dependent (eq.→area, driven→volume); PBH
evaporation tests it & REFRAME / PREDICTION &
\texttt{run\_residual\_resolution.py} \\
dS count & count = horizon Hausdorff dim; driving roughens
\(d_{\mathrm{H}}\) 2→3; residue → 2 sub-questions & ANALYSIS &
\texttt{run\_ds\_count.py} \\
\textbf{residue} & \textbf{the m≈1 volume branch is the horizon's state}
& \textbf{OPEN (L1/2) / measurable (L3)} & dS holography / DR3+Euclid \\
\end{longtable}

\subsection{Reproducibility}\label{reproducibility}

\begin{verbatim}
python reproduce_all.py --only derivA,derivB,derivC,derivD,derivE,residue,ness,jcoup,\
deposdrv,socatt,closure,litinf,violrel,socsig,v1loc,resid,bhnew,pbhdrv,dscount,kpzmem,dsmarg,dsresolve
\end{verbatim}

Figures 1--5 are generated by
\texttt{paper/make\_foundations\_figures.py}: Figs 1 (reduction chain),
2 (inferences) and 5 (SOC signature) are built directly there, while Fig
3 = \texttt{criticality\_soc.png} and Fig 4 =
\texttt{deposition\_drive.png} are produced by the
\texttt{soc}/\texttt{deposdrv} run-scripts. Each \texttt{run\_*.py}
prints a verdict and runs validation assertions; all stages PASS.

\section{Code and data availability}\label{code-and-data-availability}

All code and analysis pipelines are publicly available at
\texttt{github.com/spsingularity/SEDE-release}. Every quantitative claim
in this paper is produced by a named \texttt{run\_*.py} script wired
into the single entry point \texttt{reproduce\_all.py} (stages
\texttt{soc}, \texttt{deriv\{A–E\}}, \texttt{residue}, \texttt{ness},
\texttt{jcoup}, \texttt{deposdrv}, \texttt{socatt}, \texttt{closure},
\texttt{dscount}, \texttt{kpzmem}, \texttt{dsmarg}, \texttt{dsresolve},
\ldots), each carrying validation assertions; the figures are
regenerated by \texttt{make\_foundations\_figures.py}. The model and its
cosmology live in the \texttt{sede/} package with the companion (Paper
I); no additional observational data beyond Paper I's standard public
inputs is used here. A tagged release is archived at Zenodo, DOI
\href{https://doi.org/10.5281/zenodo.21050314}{10.5281/zenodo.21050314}.

\section{Use of AI tools}\label{use-of-ai-tools}

Artificial-intelligence tools (large language models) were used to
assist with drafting and editing the manuscript, with developing and
cross-checking the analysis code, and with literature searches. The
author conceived the model, directed and independently verified all
analyses, and takes full responsibility for the content, including all
derivations, results, and conclusions. No AI tool is an author.

\section{References}\label{references}

BibTeX in \texttt{paper/refs.bib} (cosmology, shared with Paper I) and
\texttt{paper/refs\_foundations.bib} (foundational additions for this
paper). Key entries, by theme:

\emph{Horizon thermodynamics \& entanglement} --- Jacobson,
\emph{Thermodynamics of Spacetime}, PRL 75 (1995) 1260
{[}\texttt{Jacobson:1995ab}{]}; Cai \& Kim, JHEP 02 (2005) 050
{[}\texttt{Cai:2005ra}{]}; Jacobson, \emph{Entanglement Equilibrium and
the Einstein Equation}, PRL 116 (2016) 201101, arXiv:1505.04753
{[}\texttt{Jacobson:2015hqa}{]}; Ryu \& Takayanagi, PRL 96 (2006)
181602, hep-th/0603001 {[}\texttt{Ryu:2006bv}{]}; Gibbons \& Hawking,
PRD 15 (1977) 2738 {[}\texttt{GibbonsHawking1977}{]}.

\emph{Generalised entropy \& holographic dark energy} --- Cohen, Kaplan
\& Nelson, PRL 82 (1999) 4971 {[}\texttt{Cohen:1998zx}{]}; Barrow, PLB
808 (2020) 135643 {[}\texttt{Barrow:2020tzx}{]}; Saridakis, PRD 102
(2020) 123525 {[}\texttt{Saridakis:2020zol}{]}; Tsallis \& Cirto, EPJC
73 (2013) 2487 {[}\texttt{Tsallis:2012js}{]}; Padmanabhan,
arXiv:1206.4916 {[}\texttt{Padmanabhan:2012ik}{]}; Bousso, RMP 74 (2002)
825 {[}\texttt{Bousso:2002ju}{]}.

\emph{Entropy-production acceleration \& non-equilibrium gravity} ---
García-Bellido \& Espinosa-Portalés, \emph{Cosmic acceleration from
first principles}, PDU 34 (2021) 100892, arXiv:2106.16014
{[}\texttt{Garcia-Bellido:2021idr}{]}; Espinosa-Portalés \&
García-Bellido, PDU 34 (2021) 100893, arXiv:2106.16012
{[}\texttt{Espinosa-Portales:2021cac}{]}.

\emph{Quantum chaos, ETH, de Sitter holography} --- Maldacena, Shenker
\& Stanford, JHEP 08 (2016) 106, arXiv:1503.01409
{[}\texttt{Maldacena:2015waa}{]}; Cotler et al., JHEP 05 (2017) 118
{[}\texttt{Cotler:2016fpe}{]}; Page, PRL 71 (1993) 1291
{[}\texttt{Page:1993df}{]}; Deutsch, PRA 43 (1991) 2046
{[}\texttt{Deutsch:1991msp}{]}; Srednicki, PRE 50 (1994) 888,
cond-mat/9403051 {[}\texttt{Srednicki:1994mfb}{]}; Rigol, Dunjko \&
Olshanii, Nature 452 (2008) 854 {[}\texttt{Rigol:2007juv}{]}; Banks,
IJMPA 16 (2001) 910, hep-th/0007146 {[}\texttt{Banks:2000fe}{]};
Susskind, JHAP 1 (2021) 1, arXiv:2109.14104
{[}\texttt{Susskind:2021esx}{]}; Lin, JHEP 11 (2022) 060,
arXiv:2208.07032 {[}\texttt{Lin:2022rbf}{]}; Narovlansky \& Verlinde,
JHEP 05 (2025) 032, arXiv:2310.16994 {[}\texttt{Narovlansky:2023lfz}{]}.

\emph{Emergent gravity} --- Verlinde, JHEP 04 (2011) 029
{[}\texttt{Verlinde:2010hp}{]}; Verlinde, SciPost Phys. 2 (2017) 016
{[}\texttt{Verlinde:2016toy}{]}.

\emph{Long-range systems, criticality, monitored dynamics} --- Campa,
Dauxois \& Ruffo, Phys. Rept. 480 (2009) 57, arXiv:0907.0323
{[}\texttt{Campa:2009jxa}{]}; Bak, Tang \& Wiesenfeld, PRL 59 (1987) 381
{[}\texttt{Bak:1987xua}{]}; Hawking \& Page, CMP 87 (1983) 577
{[}\texttt{Hawking:1982dh}{]}; Fisher et al., Ann. Rev.~Cond. Mat. Phys.
14 (2023) 335, arXiv:2207.14280 {[}\texttt{Fisher:2022qey}{]}.

\emph{Volume-law entanglement in cosmology} --- Khlebnikov \& Sheoran,
arXiv:1907.00487 {[}\texttt{Khlebnikov:2019wbk}{]}.

\emph{Structure-sourced dark energy} --- Gough, \emph{Information dark
energy} (≈2008) {[}\texttt{Gough2008}{]}; Pandev, \emph{Structural
Entropy Dark Energy: a fixed-parameter, growth-gated holographic
dark-energy model without a cosmological constant} (Paper I, companion
to this work; in preparation, 2025) --- the cosmology this paper's
foundations underwrite.

\end{document}
